% Theta Quantum Structure Bridge — Cross-Reference Summary
% Version: 1.0
% Date: 2026-03-10
% Status: Canonical bridge — cross-references only; no new derivations
% Author: Ing. David Jaroš
% © 2026 Ing. David Jaroš — CC BY-NC-ND 4.0
%
% PURPOSE: This file provides a single-entry navigation guide to the
% quantum structure of the Θ-field: action, canonical quantization,
% path integral, propagator, KK spectrum, and fermion emergence.
% An external reviewer should be able to assess the status of the
% Θ-field quantum program from this file alone.

\section{Quantum Structure of the \texorpdfstring{$\Theta$}{Theta}-Field: Navigation Guide}
\label{sec:bridges:theta_quantum}

\begin{tcolorbox}[colback=blue!5!white,colframe=blue!50!black,title=Θ-Field Quantum Structure Summary]
\textbf{PROVED [L0]:} Algebraic isomorphism $\mathbb{B}\cong\mathrm{Mat}(2,\mathbb{C})$ \\
\textbf{PROVED [L0]:} KK mass spectrum $m_n^2 = m^2 + n^2/R_\psi^2$ \\
\textbf{PROVED [L0]:} Three-generation structure from $\psi$-winding modes \\
\textbf{PROVED [L1]:} Canonical action $S[\Theta]$ and conjugate momenta \\
\textbf{PROVED [L1]:} Free propagator $\Delta_n(k) = i/(k^2 - m_n^2 + i\varepsilon)$ \\
\textbf{PROVED [L1]:} Dirac equation as $\psi\to 0$ limit of master equation \\
\textbf{PROVED [L1]:} Fermionic statistics and Grassmannian measure for odd-$n$ modes \\
\textbf{PROVED [L1]:} $R_\psi$ fixed via self-dual moduli potential (one-loop) \\
\textbf{[O] Open:} Yukawa matrix from moduli without fitting
\end{tcolorbox}

% -----------------------------------------------------------------------
\subsection{Step Q1: Canonical Action}
% -----------------------------------------------------------------------

\textbf{Claim:} The canonical UBT action is:
\begin{equation}
  S[\Theta]
  = \int d^4x\,d\tau\;\sqrt{|\det G|}\;
    \Bigl[
      \tfrac{1}{2}G^{MN}\langle D_M\Theta, D_N\Theta\rangle
      - V(\Theta)
      - \tfrac{1}{4}\langle F_{MN}, F^{MN}\rangle
    \Bigr].
\end{equation}

\textbf{Location (derivation):}
\texttt{consolidation\_project/appendix\_A\_theta\_action.tex} — full derivation
with boundary terms and dimensional analysis.

\textbf{Location (research):}
\texttt{research/theta\_quantum\_field.tex} §1 — canonical form and sources.

\textbf{Status: [L1]} — proved given axioms.

% -----------------------------------------------------------------------
\subsection{Step Q2: Canonical Momenta and Commutation Relations}
% -----------------------------------------------------------------------

\textbf{Claim:} The canonical momentum is $\Pi^a = \dot\Theta_a$, and the
equal-time commutator is:
\begin{equation}
  [\Theta_a(x,\psi), \Pi^b(x',\psi')] = i\hbar\,\delta^{ab}\,
  \delta^{(3)}(x-x')\,\delta(\psi-\psi').
\end{equation}

\textbf{Location:} \texttt{research/theta\_quantum\_field.tex} §2.

\textbf{Status: [L1]}.

% -----------------------------------------------------------------------
\subsection{Step Q3: Kaluza--Klein Spectrum}
% -----------------------------------------------------------------------

\textbf{Claim:} The KK expansion
$\Theta(x,t,\psi) = \sum_n \hat\Theta_n(x,t)\,e^{in\psi/R_\psi}$
yields effective masses $m_n^2 = m^2 + n^2/R_\psi^2$.

\textbf{Location (primary):} \texttt{report/theta\_spectrum\_analysis.md} §2.2;
\texttt{canonical/geometry/phase\_projection.tex} §KK.

\textbf{Location (research):} \texttt{research/theta\_quantum\_field.tex} §4.

\textbf{Status: [L0]}.

% -----------------------------------------------------------------------
\subsection{Step Q4: Free Propagator}
% -----------------------------------------------------------------------

\textbf{Claim:} The Feynman propagator for KK mode $n$ in momentum space is:
\begin{equation}
  \Delta_n(k) = \frac{i}{k^\mu k_\mu - m_n^2 + i\varepsilon}.
\end{equation}

\textbf{Location:} \texttt{research/theta\_quantum\_field.tex} §3.2,
Definition~3.2.

\textbf{Status: [L1]}.

% -----------------------------------------------------------------------
\subsection{Step Q5: Dirac Equation as Effective Limit}
% -----------------------------------------------------------------------

\textbf{Claim:} In the $\psi\to 0$ limit with $\Theta\in\mathcal{S}$
(spinor subspace), the master equation reduces to the Dirac equation
$(i\gamma^\mu D_\mu - m)\psi = 0$.

\textbf{Location (full proof):}
\texttt{canonical/bridges/QED\_complete\_chain.tex} Theorem D.4;
\texttt{consolidation\_project/FPE\_verification/step3\_dirac\_emergence.tex}.

\textbf{Location (research):} \texttt{research/theta\_fermion\_emergence.tex} §2.

\textbf{Status: [L1]}.

% -----------------------------------------------------------------------
\subsection{Step Q6: Fermionic Statistics}
% -----------------------------------------------------------------------

\textbf{Claim:} Odd-winding KK modes ($n$ odd) carry $\mathcal{P}_\psi = -1$
parity, and the path-integral measure for these modes is Grassmannian,
yielding canonical anticommutation relations.

\textbf{Location (primary):} \texttt{canonical/bridges/fermionic\_statistics\_bridge.tex}
— exchange-phase holonomy proof with canonical anticommutation relations.

\textbf{Location (background):} \texttt{research/theta\_fermion\_emergence.tex} §3.

\textbf{Status: [L1]} — proved via exchange-phase argument
(see \texttt{canonical/bridges/fermionic\_statistics\_bridge.tex}
Theorems~2 and~3).

\begin{tcolorbox}[colback=green!5!white,colframe=green!50!black,title=Gap Q6 Closed]
  The Grassmannian path-integral measure for odd-winding KK modes has been
  derived in \texttt{canonical/bridges/fermionic\_statistics\_bridge.tex}.
  The proof uses the exchange-phase holonomy: adiabatic exchange of two
  $n$-th KK mode excitations in $\mathbb{R}^4\times S^1_\psi$ acquires the
  phase $(-1)^n$, forcing Grassmann statistics for $n$ odd.
  Canonical anticommutation relations follow at level \textbf{[L1]}.
\end{tcolorbox}

% -----------------------------------------------------------------------
\subsection{Step Q7: Three-Generation Structure (Proved)}
% -----------------------------------------------------------------------

\textbf{Claim:} Modes $n=1,2,3$ provide three independent fermion families.

\textbf{Location:} \texttt{DERIVATION\_INDEX.md}: ``$\psi$-modes as independent
B-fields [L0]: Proven''; \texttt{research/theta\_fermion\_emergence.tex}
Proposition~3.2.

\textbf{Status: [L0]}.

% -----------------------------------------------------------------------
\subsection{Connection to \texorpdfstring{$\alpha$}{alpha} and \texorpdfstring{$B_{\mathrm{base}}$}{B\_base}}
% -----------------------------------------------------------------------

The KK mode sum over all $n$ contributes to the vacuum polarization and
ultimately to the $B_{\mathrm{base}}$ coefficient in the fine-structure
constant formula.  The complete derivation $B_{\mathrm{base}} = N_{\mathrm{eff}}^{3/2}
= 12^{3/2} \approx 41.57$ is in:
\texttt{canonical/interactions/B\_base\_derivation\_complete.tex}.
Cross-references:
\texttt{canonical/bridges/QED\_limit\_bridge.tex},
\texttt{research/B\_base\_spectral\_determinant.tex}.

\end{document}
